\chapter{Thảo luận}
\label{chap:discussion}

\section{Kết quả đạt được}

Đồ án đã đạt được một nền tảng tương đối đầy đủ để trình bày DevSecOps cho web microservices. Về ứng dụng, SageLMS có frontend, gateway, nhiều backend services, OpenAPI contracts, database migrations và Dockerfile. Về pipeline, repo có PR validation, CD build/scan/push/deploy PR, infra plan và post-deploy check. Về cloud runtime, nhóm đã có GKE, CloudNativePG, Redis, Harbor, ESO và các manifests GitOps/Kustomize.

Điểm quan trọng là nhóm không chỉ mô tả pipeline trên lý thuyết mà đã có evidence: Trivy reports, SBOM, image digests, deploy summary, OpenTofu outputs và cloud foundation evidence.

\section{Điểm mạnh}

Thiết kế pipeline có một số điểm mạnh:

\begin{itemize}
  \item Vai trò công cụ rõ ràng, tránh dùng nhiều scanner trùng mục đích.
  \item CI không deploy trực tiếp vào cluster, giúp tách artifact creation khỏi runtime reconciliation.
  \item Image được gắn với git SHA và digest, tăng khả năng truy vết.
  \item SBOM và signing đưa supply chain security vào báo cáo thay vì chỉ dừng ở Docker build.
  \item OpenTofu và Checkov giúp hạ tầng có thể review và kiểm soát.
  \item Shared Cloud DevSecOps Environment phù hợp với nguồn lực sinh viên nhưng vẫn chứng minh được luồng production-oriented.
\end{itemize}

\section{Hạn chế}

Một số hạn chế cần trình bày trung thực:

\begin{itemize}
  \item Đồ án chưa tách đủ ba môi trường dev, staging và production như hệ thống thương mại.
  \item Restore drill CloudNativePG cần được hoàn thiện nếu muốn chứng minh disaster recovery đầy đủ.
  \item Observability có thể mới ở mức định hướng hoặc foundation, chưa chắc đã có dashboard/alert đầy đủ.
  \item Kyverno policy có thể cần chuyển dần từ audit sang enforce để tránh làm gián đoạn demo.
  \item AI Tutor và worker nên được xác định rõ là nằm trong demo runtime chính hay nằm ở phạm vi mở rộng.
  \item Một số thông tin như MSSV, giảng viên, ảnh chụp workflow và kết quả demo cuối cần bổ sung trước khi nộp.
\end{itemize}

\section{So sánh với production thực tế}

Nếu triển khai production thật, pipeline cần được mở rộng thêm. Hệ thống nên có ít nhất staging và production riêng, policy enforce chặt hơn, secret rotation, backup restore drill định kỳ, SLO/SLA, monitoring dashboard, alerting rule, incident response runbook và progressive delivery như canary hoặc blue-green. Ngoài ra, supply chain có thể nâng cấp theo hướng SLSA provenance, keyless signing chuẩn hóa bằng OIDC và admission policy xác minh chữ ký image trước khi pod chạy.

Tuy nhiên, với phạm vi đồ án môn học, việc duy trì một môi trường GKE dùng chung là lựa chọn hợp lý. Nó giảm chi phí nhưng vẫn cho phép chứng minh phần cốt lõi của DevSecOps: kiểm tra tự động, artifact đáng tin cậy, GitOps deployment, IaC và evidence vận hành.

\section{Hướng phát triển}

Các hướng phát triển tiếp theo gồm:

\begin{itemize}
  \item Hoàn thiện restore drill CloudNativePG sang cluster tạm.
  \item Bật Kyverno verifyImages để enforce image từ Harbor và Cosign signature.
  \item Hoàn thiện dashboard Prometheus/Grafana/Loki/Tempo cho các service.
  \item Bổ sung e2e test cho demo scenario chính.
  \item Chuẩn hóa Cosign keyless signing bằng GitHub OIDC.
  \item Tách staging/production nếu có tài nguyên cloud.
  \item Tự động hóa release notes và evidence packaging cho mỗi lần deploy.
\end{itemize}
